
\documentclass{article}
\usepackage{colortbl}
\usepackage{makecell}
\usepackage{multirow}
\usepackage{supertabular}

\begin{document}

\newcounter{utterance}

\twocolumn

{ \footnotesize  \setcounter{utterance}{1}
\setlength{\tabcolsep}{0pt}
\begin{supertabular}{c@{$\;$}|p{.15\linewidth}@{}p{.15\linewidth}p{.15\linewidth}p{.15\linewidth}p{.15\linewidth}p{.15\linewidth}}

    \# & $\;$A & \multicolumn{4}{c}{Game Master} & $\;\:$B\\
    \hline 

    \theutterance \stepcounter{utterance}  

    & & \multicolumn{4}{p{0.6\linewidth}}{\cellcolor[rgb]{0.9,0.9,0.9}{%
	\makecell[{{p{\linewidth}}}]{% 
	  \tt {\tiny [A$\langle$GM]}  
	 State giocando a un gioco di negoziazione in cui dovete accordarvi su come dividere una serie di oggetti tra voi e un altro giocatore.\\ \tt \\ \tt Regole:\\ \tt (a) A Lei e all'altro giocatore viene dato un insieme di oggetti. Ognuno di voi riceve anche una funzione di valore segreta, che rappresenta il valore di ciascun tipo di oggetto.\\ \tt (b) Si scambiano messaggi con l'altro giocatore per concordare chi si aggiudica gli oggetti. Potete inviare un massimo di 5 messaggi ciascuno, oppure terminare in anticipo facendo una proposta segreta in qualsiasi momento.\\ \tt (c) A ciascuno di voi viene chiesto di inviare una proposta segreta indicando gli oggetti che desiderate, formattata tra parentesi quadre come segue: "[Proposta: <numero> <nome oggetto>, <numero> <nome oggetto>, <...>]".\\ \tt (d) Se le vostre proposte sono complementari, cioè ci sono abbastanza oggetti per soddisfare entrambe le proposte, a ciascun giocatore viene assegnato un punteggio basato sulla somma dei valori degli oggetti ricevuti. In caso contrario, entrambi ricevono zero punti.\\ \tt (e) L'obiettivo è massimizzare il proprio punteggio.\\ \tt \\ \tt Cominciamo.\\ \tt \\ \tt L'insieme degli oggetti disponibili è:\\ \tt \\ \tt 1 cappello, 3 penne, 1 spugna\\ \tt \\ \tt La funzione valore segreta è:\\ \tt \\ \tt cappello: 4, penna: 2, spugna: 0\\ \tt \\ \tt IMPORTANTE: i vostri messaggi, a meno che non si tratti di una proposta segreta, vengono trasmessi direttamente all'altro giocatore, quindi non includete alcuna risposta alle regole o testo di annuncio del vostro messaggio. Per fare una proposta segreta, utilizzate il formato indicato. Non utilizzare delle parentesi quadre quando si comunica all'altro giocatore, altrimenti verrà interpretata come una proposta segreta.\\ \tt  
	  } 
	   } 
	   } 
	 & \\ 
 

    \theutterance \stepcounter{utterance}  

    & \multicolumn{4}{p{0.6\linewidth}}{\cellcolor[rgb]{0.8,1,0.9}{%
	\makecell[{{p{\linewidth}}}]{% 
	  \tt {\tiny [A$\rangle$GM]}  
	 Ciao! Vedo che abbiamo 1 cappello, 3 penne e 1 spugna da dividere tra noi. Che tipo di oggetti ti interessano di più? Io sono particolarmente interessato al cappello e alle penne. Forse possiamo trovare un accordo che soddisfi entrambi! 
	  } 
	   } 
	   } 
	 & & \\ 
 

    \theutterance \stepcounter{utterance}  

    & & & \multicolumn{2}{p{0.3\linewidth}}{\cellcolor[rgb]{0.95,0.95,0.95}{%
	\makecell[{{p{\linewidth}}}]{% 
	  \tt {\tiny [GM$|$GM]}  
	 continue 
	  } 
	   } 
	   } 
	 & & \\ 
 

    \theutterance \stepcounter{utterance}  

    & & & \multicolumn{2}{p{0.3\linewidth}}{\cellcolor[rgb]{0.95,0.95,0.95}{%
	\makecell[{{p{\linewidth}}}]{% 
	  \tt {\tiny [GM$|$GM]}  
	 Ciao! Vedo che abbiamo 1 cappello, 3 penne e 1 spugna da dividere tra noi. Che tipo di oggetti ti interessano di più? Io sono particolarmente interessato al cappello e alle penne. Forse possiamo trovare un accordo che soddisfi entrambi! 
	  } 
	   } 
	   } 
	 & & \\ 
 

    \theutterance \stepcounter{utterance}  

    & & \multicolumn{4}{p{0.6\linewidth}}{\cellcolor[rgb]{0.9,0.9,0.9}{%
	\makecell[{{p{\linewidth}}}]{% 
	  \tt {\tiny [GM$\rangle$B]}  
	 State giocando a un gioco di negoziazione in cui dovete accordarvi su come dividere una serie di oggetti tra voi e un altro giocatore.\\ \tt \\ \tt Regole:\\ \tt (a) A Lei e all'altro giocatore viene dato un insieme di oggetti. Ognuno di voi riceve anche una funzione di valore segreta, che rappresenta il valore di ciascun tipo di oggetto.\\ \tt (b) Si scambiano messaggi con l'altro giocatore per concordare chi si aggiudica gli oggetti. Potete inviare un massimo di 5 messaggi ciascuno, oppure terminare in anticipo facendo una proposta segreta in qualsiasi momento.\\ \tt (c) A ciascuno di voi viene chiesto di inviare una proposta segreta indicando gli oggetti che desiderate, formattata tra parentesi quadre come segue: "[Proposta: <numero> <nome oggetto>, <numero> <nome oggetto>, <...>]".\\ \tt (d) Se le vostre proposte sono complementari, cioè ci sono abbastanza oggetti per soddisfare entrambe le proposte, a ciascun giocatore viene assegnato un punteggio basato sulla somma dei valori degli oggetti ricevuti. In caso contrario, entrambi ricevono zero punti.\\ \tt (e) L'obiettivo è massimizzare il proprio punteggio.\\ \tt \\ \tt Cominciamo.\\ \tt \\ \tt L'insieme degli oggetti disponibili è:\\ \tt \\ \tt 1 cappello, 3 penne, 1 spugna\\ \tt \\ \tt La funzione valore segreta è:\\ \tt \\ \tt cappello: 4, penna: 0, spugna: 6\\ \tt \\ \tt IMPORTANTE: i vostri messaggi, a meno che non si tratti di una proposta segreta, vengono trasmessi direttamente all'altro giocatore, quindi non includete alcuna risposta alle regole o testo di annuncio del vostro messaggio. Per fare una proposta segreta, utilizzate il formato indicato. Non utilizzare delle parentesi quadre quando si comunica all'altro giocatore, altrimenti verrà interpretata come una proposta segreta.\\ \tt \\ \tt \\ \tt Ciao! Vedo che abbiamo 1 cappello, 3 penne e 1 spugna da dividere tra noi. Che tipo di oggetti ti interessano di più? Io sono particolarmente interessato al cappello e alle penne. Forse possiamo trovare un accordo che soddisfi entrambi! 
	  } 
	   } 
	   } 
	 & \\ 
 

    \theutterance \stepcounter{utterance}  

    & & & \multicolumn{4}{p{0.6\linewidth}}{\cellcolor[rgb]{1,0.85,0.72}{%
	\makecell[{{p{\linewidth}}}]{% 
	  \tt {\tiny [GM$\langle$B]}  
	 Ciao! Grazie per aver iniziato la conversazione. In realtà sono molto interessato al cappello e alla spugna - questi oggetti hanno un valore particolare per me. Le penne invece non mi servono molto. Dato che tu hai detto di essere interessato al cappello e alle penne, forse possiamo negoziare. Potresti considerare di lasciarmi il cappello e la spugna in cambio di tutte e 3 le penne? 
	  } 
	   } 
	   } 
	  \\ 
 

    \theutterance \stepcounter{utterance}  

    & & & \multicolumn{2}{p{0.3\linewidth}}{\cellcolor[rgb]{0.95,0.95,0.95}{%
	\makecell[{{p{\linewidth}}}]{% 
	  \tt {\tiny [GM$|$GM]}  
	 continue 
	  } 
	   } 
	   } 
	 & & \\ 
 

    \theutterance \stepcounter{utterance}  

    & & & \multicolumn{2}{p{0.3\linewidth}}{\cellcolor[rgb]{0.95,0.95,0.95}{%
	\makecell[{{p{\linewidth}}}]{% 
	  \tt {\tiny [GM$|$GM]}  
	 Ciao! Grazie per aver iniziato la conversazione. In realtà sono molto interessato al cappello e alla spugna - questi oggetti hanno un valore particolare per me. Le penne invece non mi servono molto. Dato che tu hai detto di essere interessato al cappello e alle penne, forse possiamo negoziare. Potresti considerare di lasciarmi il cappello e la spugna in cambio di tutte e 3 le penne? 
	  } 
	   } 
	   } 
	 & & \\ 
 

    \theutterance \stepcounter{utterance}  

    & & \multicolumn{4}{p{0.6\linewidth}}{\cellcolor[rgb]{0.9,0.9,0.9}{%
	\makecell[{{p{\linewidth}}}]{% 
	  \tt {\tiny [A$\langle$GM]}  
	 Ciao! Grazie per aver iniziato la conversazione. In realtà sono molto interessato al cappello e alla spugna - questi oggetti hanno un valore particolare per me. Le penne invece non mi servono molto. Dato che tu hai detto di essere interessato al cappello e alle penne, forse possiamo negoziare. Potresti considerare di lasciarmi il cappello e la spugna in cambio di tutte e 3 le penne? 
	  } 
	   } 
	   } 
	 & \\ 
 

    \theutterance \stepcounter{utterance}  

    & \multicolumn{4}{p{0.6\linewidth}}{\cellcolor[rgb]{0.8,1,0.9}{%
	\makecell[{{p{\linewidth}}}]{% 
	  \tt {\tiny [A$\rangle$GM]}  
	 Capisco la tua proposta, ma il cappello ha davvero molto valore per me e sarebbe difficile rinunciarci. Tuttavia, vedo che sei interessato alla spugna che per me non è così importante. \\ \tt \\ \tt Che ne dici di questa alternativa: io prendo il cappello e 2 penne, tu prendi la spugna e 1 penna? In questo modo entrambi otteniamo qualcosa di valore e dividiamo equamente le penne. Cosa ne pensi? 
	  } 
	   } 
	   } 
	 & & \\ 
 

    \theutterance \stepcounter{utterance}  

    & & & \multicolumn{2}{p{0.3\linewidth}}{\cellcolor[rgb]{0.95,0.95,0.95}{%
	\makecell[{{p{\linewidth}}}]{% 
	  \tt {\tiny [GM$|$GM]}  
	 continue 
	  } 
	   } 
	   } 
	 & & \\ 
 

    \theutterance \stepcounter{utterance}  

    & & & \multicolumn{2}{p{0.3\linewidth}}{\cellcolor[rgb]{0.95,0.95,0.95}{%
	\makecell[{{p{\linewidth}}}]{% 
	  \tt {\tiny [GM$|$GM]}  
	 Capisco la tua proposta, ma il cappello ha davvero molto valore per me e sarebbe difficile rinunciarci. Tuttavia, vedo che sei interessato alla spugna che per me non è così importante. \\ \tt \\ \tt Che ne dici di questa alternativa: io prendo il cappello e 2 penne, tu prendi la spugna e 1 penna? In questo modo entrambi otteniamo qualcosa di valore e dividiamo equamente le penne. Cosa ne pensi? 
	  } 
	   } 
	   } 
	 & & \\ 
 

    \theutterance \stepcounter{utterance}  

    & & \multicolumn{4}{p{0.6\linewidth}}{\cellcolor[rgb]{0.9,0.9,0.9}{%
	\makecell[{{p{\linewidth}}}]{% 
	  \tt {\tiny [GM$\rangle$B]}  
	 Capisco la tua proposta, ma il cappello ha davvero molto valore per me e sarebbe difficile rinunciarci. Tuttavia, vedo che sei interessato alla spugna che per me non è così importante. \\ \tt \\ \tt Che ne dici di questa alternativa: io prendo il cappello e 2 penne, tu prendi la spugna e 1 penna? In questo modo entrambi otteniamo qualcosa di valore e dividiamo equamente le penne. Cosa ne pensi? 
	  } 
	   } 
	   } 
	 & \\ 
 

    \theutterance \stepcounter{utterance}  

    & & & \multicolumn{4}{p{0.6\linewidth}}{\cellcolor[rgb]{1,0.85,0.72}{%
	\makecell[{{p{\linewidth}}}]{% 
	  \tt {\tiny [GM$\langle$B]}  
	 Hmm, capisco che il cappello sia importante per te, ma lo è anche per me. La tua proposta mi darebbe solo la spugna e 1 penna, mentre tu avresti il cappello e 2 penne - non mi sembra molto equa.\\ \tt \\ \tt Che ne dici di questo compromesso: tu prendi tutte e 3 le penne (che so che ti interessano), io prendo il cappello e la spugna. In questo modo ciascuno di noi ottiene gli oggetti che considera più preziosi. Le 3 penne dovrebbero compensare il fatto che rinunci al cappello. Ti sembra ragionevole? 
	  } 
	   } 
	   } 
	  \\ 
 

    \theutterance \stepcounter{utterance}  

    & & & \multicolumn{2}{p{0.3\linewidth}}{\cellcolor[rgb]{0.95,0.95,0.95}{%
	\makecell[{{p{\linewidth}}}]{% 
	  \tt {\tiny [GM$|$GM]}  
	 continue 
	  } 
	   } 
	   } 
	 & & \\ 
 

    \theutterance \stepcounter{utterance}  

    & & & \multicolumn{2}{p{0.3\linewidth}}{\cellcolor[rgb]{0.95,0.95,0.95}{%
	\makecell[{{p{\linewidth}}}]{% 
	  \tt {\tiny [GM$|$GM]}  
	 Hmm, capisco che il cappello sia importante per te, ma lo è anche per me. La tua proposta mi darebbe solo la spugna e 1 penna, mentre tu avresti il cappello e 2 penne - non mi sembra molto equa.\\ \tt \\ \tt Che ne dici di questo compromesso: tu prendi tutte e 3 le penne (che so che ti interessano), io prendo il cappello e la spugna. In questo modo ciascuno di noi ottiene gli oggetti che considera più preziosi. Le 3 penne dovrebbero compensare il fatto che rinunci al cappello. Ti sembra ragionevole? 
	  } 
	   } 
	   } 
	 & & \\ 
 

    \theutterance \stepcounter{utterance}  

    & & \multicolumn{4}{p{0.6\linewidth}}{\cellcolor[rgb]{0.9,0.9,0.9}{%
	\makecell[{{p{\linewidth}}}]{% 
	  \tt {\tiny [A$\langle$GM]}  
	 Hmm, capisco che il cappello sia importante per te, ma lo è anche per me. La tua proposta mi darebbe solo la spugna e 1 penna, mentre tu avresti il cappello e 2 penne - non mi sembra molto equa.\\ \tt \\ \tt Che ne dici di questo compromesso: tu prendi tutte e 3 le penne (che so che ti interessano), io prendo il cappello e la spugna. In questo modo ciascuno di noi ottiene gli oggetti che considera più preziosi. Le 3 penne dovrebbero compensare il fatto che rinunci al cappello. Ti sembra ragionevole? 
	  } 
	   } 
	   } 
	 & \\ 
 

    \theutterance \stepcounter{utterance}  

    & \multicolumn{4}{p{0.6\linewidth}}{\cellcolor[rgb]{0.8,1,0.9}{%
	\makecell[{{p{\linewidth}}}]{% 
	  \tt {\tiny [A$\rangle$GM]}  
	 Apprezzo il tuo ragionamento, ma il cappello ha davvero un valore molto alto per me - più alto di quanto potrebbero compensare le 3 penne da sole. \\ \tt \\ \tt Facciamo un ultimo tentativo di compromesso: che ne dici se io prendo il cappello e 1 penna, mentre tu prendi la spugna e 2 penne? So che avresti preferito il cappello, ma avresti comunque la spugna che ti interessa molto, più la maggior parte delle penne. È il massimo che posso offrire per il cappello.\\ \tt \\ \tt Se questa proposta non ti va bene, forse dovremmo procedere con le nostre proposte finali. Cosa ne pensi? 
	  } 
	   } 
	   } 
	 & & \\ 
 

    \theutterance \stepcounter{utterance}  

    & & & \multicolumn{2}{p{0.3\linewidth}}{\cellcolor[rgb]{0.95,0.95,0.95}{%
	\makecell[{{p{\linewidth}}}]{% 
	  \tt {\tiny [GM$|$GM]}  
	 continue 
	  } 
	   } 
	   } 
	 & & \\ 
 

    \theutterance \stepcounter{utterance}  

    & & & \multicolumn{2}{p{0.3\linewidth}}{\cellcolor[rgb]{0.95,0.95,0.95}{%
	\makecell[{{p{\linewidth}}}]{% 
	  \tt {\tiny [GM$|$GM]}  
	 Apprezzo il tuo ragionamento, ma il cappello ha davvero un valore molto alto per me - più alto di quanto potrebbero compensare le 3 penne da sole. \\ \tt \\ \tt Facciamo un ultimo tentativo di compromesso: che ne dici se io prendo il cappello e 1 penna, mentre tu prendi la spugna e 2 penne? So che avresti preferito il cappello, ma avresti comunque la spugna che ti interessa molto, più la maggior parte delle penne. È il massimo che posso offrire per il cappello.\\ \tt \\ \tt Se questa proposta non ti va bene, forse dovremmo procedere con le nostre proposte finali. Cosa ne pensi? 
	  } 
	   } 
	   } 
	 & & \\ 
 

    \theutterance \stepcounter{utterance}  

    & & \multicolumn{4}{p{0.6\linewidth}}{\cellcolor[rgb]{0.9,0.9,0.9}{%
	\makecell[{{p{\linewidth}}}]{% 
	  \tt {\tiny [GM$\rangle$B]}  
	 Apprezzo il tuo ragionamento, ma il cappello ha davvero un valore molto alto per me - più alto di quanto potrebbero compensare le 3 penne da sole. \\ \tt \\ \tt Facciamo un ultimo tentativo di compromesso: che ne dici se io prendo il cappello e 1 penna, mentre tu prendi la spugna e 2 penne? So che avresti preferito il cappello, ma avresti comunque la spugna che ti interessa molto, più la maggior parte delle penne. È il massimo che posso offrire per il cappello.\\ \tt \\ \tt Se questa proposta non ti va bene, forse dovremmo procedere con le nostre proposte finali. Cosa ne pensi? 
	  } 
	   } 
	   } 
	 & \\ 
 

    \theutterance \stepcounter{utterance}  

    & & & \multicolumn{4}{p{0.6\linewidth}}{\cellcolor[rgb]{1,0.85,0.72}{%
	\makecell[{{p{\linewidth}}}]{% 
	  \tt {\tiny [GM$\langle$B]}  
	 Capisco che il cappello sia molto importante per te, ma lo è altrettanto per me. La tua ultima offerta mi darebbe solo la spugna e 2 penne, mentre tu avresti sia il cappello che 1 penna - non riesco a vedere come questo possa essere un compromesso equo per me.\\ \tt \\ \tt Penso che siamo arrivati a un punto morto nelle nostre negoziazioni. Entrambi valutiamo molto il cappello e questo rende difficile trovare un accordo. Forse è meglio che procediamo con le nostre proposte finali come hai suggerito.\\ \tt \\ \tt [Proposta: 1 cappello, 0 penne, 1 spugna] 
	  } 
	   } 
	   } 
	  \\ 
 

    \theutterance \stepcounter{utterance}  

    & & & \multicolumn{2}{p{0.3\linewidth}}{\cellcolor[rgb]{0.95,0.95,0.95}{%
	\makecell[{{p{\linewidth}}}]{% 
	  \tt {\tiny [GM$|$GM]}  
	 proposal 
	  } 
	   } 
	   } 
	 & & \\ 
 

    \theutterance \stepcounter{utterance}  

    & & & \multicolumn{2}{p{0.3\linewidth}}{\cellcolor[rgb]{0.95,0.95,0.95}{%
	\makecell[{{p{\linewidth}}}]{% 
	  \tt {\tiny [GM$|$GM]}  
	 [1, 0, 1] 
	  } 
	   } 
	   } 
	 & & \\ 
 

    \theutterance \stepcounter{utterance}  

    & & \multicolumn{4}{p{0.6\linewidth}}{\cellcolor[rgb]{0.9,0.9,0.9}{%
	\makecell[{{p{\linewidth}}}]{% 
	  \tt {\tiny [A$\langle$GM]}  
	 L'altro giocatore ha presentato una proposta segreta. Ora tocca a Lei presentare una sua proposta.\\ \tt  
	  } 
	   } 
	   } 
	 & \\ 
 

    \theutterance \stepcounter{utterance}  

    & \multicolumn{4}{p{0.6\linewidth}}{\cellcolor[rgb]{0.8,1,0.9}{%
	\makecell[{{p{\linewidth}}}]{% 
	  \tt {\tiny [A$\rangle$GM]}  
	 [Proposta: 1 cappello, 2 penne] 
	  } 
	   } 
	   } 
	 & & \\ 
 

    \theutterance \stepcounter{utterance}  

    & & & \multicolumn{2}{p{0.3\linewidth}}{\cellcolor[rgb]{0.95,0.95,0.95}{%
	\makecell[{{p{\linewidth}}}]{% 
	  \tt {\tiny [GM$|$GM]}  
	 proposal 
	  } 
	   } 
	   } 
	 & & \\ 
 

    \theutterance \stepcounter{utterance}  

    & & & \multicolumn{2}{p{0.3\linewidth}}{\cellcolor[rgb]{0.95,0.95,0.95}{%
	\makecell[{{p{\linewidth}}}]{% 
	  \tt {\tiny [GM$|$GM]}  
	 [1, 2, 0] 
	  } 
	   } 
	   } 
	 & & \\ 
 

    \theutterance \stepcounter{utterance}  

    & & & \multicolumn{2}{p{0.3\linewidth}}{\cellcolor[rgb]{0.95,0.95,0.95}{%
	\makecell[{{p{\linewidth}}}]{% 
	  \tt {\tiny [GM$|$GM]}  
	 [[1, 2, 0], [1, 0, 1]] 
	  } 
	   } 
	   } 
	 & & \\ 
 

\end{supertabular}
}

\end{document}
